\section{Precision against the Trace}\label{sec:precision}

The rules in Sect.~\ref{sec:semantics} are checked against the
program's actual execution, not against a join over the program's
possible executions. This is the discipline's main precision
advantage over the corresponding static check, and the reason we
expect a runtime model in this setting to be useful at all.

\paragraph{What the runtime knows that static analysis does not.}
The simplest illustration is a program with a conditional move:
\begin{lstlisting}
let x = ...;
let r = &x.f;
if rare_condition() { move x }   // invalidates subtree(x) only here
*r                               // safe whenever rare_condition() is false
\end{lstlisting}
A sound static analysis sees the move on one branch and has to
conclude that \texttt{r} may be invalid by the time \texttt{*r} runs;
it joins the two branches and reports the union. The discipline
sees something more specific. By the time control reaches \texttt{*r},
\texttt{rare\_condition()} has either evaluated to true or to false,
and the runtime knows which. On the branch where the move did not
happen, \texttt{r} is still attached to its node, no event has
poisoned it, and the read succeeds. The discipline accepts a
strictly larger set of executions than the static check, by
construction.

\paragraph{Relation to static checking.}
We are not trying to port a static borrow checker to runtime. The
rules in Sect.~\ref{sec:semantics} were calibrated so that each rule
poisons a reference only on cases that the corresponding static
analysis would also reject when reached on every branch; this is the
sense in which the discipline is meant to be at least as permissive
as its static counterpart along a concrete trace. We do not prove
that here. The natural proof obligation, for each static rule, is to
identify a runtime invariant strong enough to rule out the same
unsafe events. This is the most important open piece of the work,
discussed further in Sect.~\ref{sec:discussion}.

\paragraph{Where the cost is worth it.}
Per-instruction shadow-state updates are not free. The discipline
pays its way in contexts where static checking is unavailable or
untrusted---bytecode that bypasses the verifier (fuzzing harnesses,
mutated bytecode, dynamic-dispatch boundaries), and effects whose
body is opaque to the bytecode language (native code). On
verifier-protected execution paths, the discipline can be turned off
entirely, or shifted into a trace-and-replay mode that records the
relevant events during execution and checks them offline.
